\documentclass[spanish, 12pt]{thesis}
\usepackage[T1]{fontenc}
\usepackage[utf8]{inputenc}
\usepackage{lmodern}
\usepackage[a4paper]{geometry}
\usepackage[spanish]{babel}
\usepackage{hyperref}
\usepackage{csquotes}

% Bibliorafía
\usepackage[style=apa,sorting=nyt,backend=biber]{biblatex}
\DeclareLanguageMapping{spanish}{spanish-apa}
\bibliography{references}
%

%\usepackage[]{parskip}
\usepackage{multicol}
\usepackage{multirow}
\usepackage[x11names,table]{xcolor} 
\usepackage{floatrow}
\usepackage{graphicx}
\usepackage{amsmath, latexsym, amsfonts, amssymb} 
\usepackage{booktabs}
\usepackage{url}
%\usepackage{float}
\usepackage[shortlabels]{enumitem}
\usepackage{tikz}
\usepackage{floatrow}
\usepackage[bf,labelsep=period]{caption}
\captionsetup{justification = centering}
\renewcommand{\baselinestretch}{1.5} % interlineado
\usepackage[bottom]{footmisc}

\usepackage{mathtools}
\usepackage{optidef}

% Comentarios
%\usepackage{comment}
%\usepackage{todonotes}
%\usepackage{easyReview}
%

% Python
\usepackage{listings}
%\usepackage{listingsutf8}
\usepackage{color}

\definecolor{mygreen}{rgb}{0,0.6,0}
\definecolor{mygray}{rgb}{0.5,0.5,0.5}
\definecolor{mymauve}{rgb}{0.58,0,0.82}
\definecolor{myred}{rgb}{0.58,0,0}


\lstset{ 
	backgroundcolor=\color{white},
	basicstyle=\footnotesize,
	breakatwhitespace=false,  
	breaklines=true,         
	captionpos=b,            
	commentstyle=\color{mygray},   
	deletekeywords={...},   
	escapeinside={\%*}{*)},         
	extendedchars=true,             
	firstnumber=1,               
	frame=,	     
%	identifierstyle=\color{blue},             
	keepspaces=true,                
	keywordstyle=\bfseries\color{mygreen},      
	language=Python,                
	morekeywords={*, ..., Var, AbstractModel, Param, Set, Obective, Constraint, Complementarity},            
	numbers=none,   % Posición de los números de línea (none, left, right).                 
	numbersep=5pt,                   
	numberstyle=\tiny\color{mygray},
	rulecolor=\color{blue},        
	showspaces=false,                
	showstringspaces=false,          
	showtabs=false,                
	stepnumber=1,                    
	stringstyle=\color{myred},     
	tabsize=2,	                  
	title=\lstname                  
}
%

\newcommand{\R}{\mathbb{R}} 
\newcommand{\N}{\mathbb{N}} 
\newcommand{\Z}{\mathbb{Z}} 
\newcommand{\Lagr}{\mathcal{L}}


\begin{document}
	\frontmatter
	\begin{titlepage}
	
	\title{\vspace*{-70pt}
		\begin{figure}
			\includegraphics[scale=0.2]{img/logo.png}
		\end{figure} 
		\vspace*{10pt}Facultad de Matemática y Computación \\
		Universidad de La Habana \\
		\vspace*{20pt}	
		Tesis de diploma de la especialidad \\ 
		Ciencia de la Computación \\
		
		\vspace*{20pt}
		Un Estudio sobre
		Problema Prueba para Evaluar la Calidad de los Algortimos %\footnote{Proyecto Nacional de Ciencias Básicas: PN223H010-005 Desarrollo de modelos y métodos matemáticos para la toma de decisiones.}
		\vspace*{-35pt}
		\author{{\large Autor: Francisco Vicente Suárez Bellón \vspace*{15pt}} \\
			{\large Tutora: Dr. C. Gemayqzel Bouza Allende} \\ \hspace{10pt}  }
		
		\date{La Habana, \today}
	}
	
\end{titlepage}

	\maketitle
	
	\chapter*{Resumen}

El problema de optimización binivel se define como minimizar una función sobre un conjunto determinado por los puntos óptimos de un modelo de programación matemática. La optimización en el nivel inferior depende de las decisiones tomadas en el nivel superior, creando así una relación de interdependencia entre ambos niveles.

Para abordar este problema, se considera la creación de un problema relacionado en la cual el nivel inferior se sustituye por las condiciones necesarias de optimalidad, siendo este un problema con restricciones de complementariedad.

En este trabajo se propone una forma de generar problemas de dos niveles cuyo problema con restricciones de complementariedad (MPEC) tiene un punto estacionario perteneciente a una de las siguientes 
clases: fuertemente estacionario, M-estacionario o C-estacionario, dependiendo de los multiplicadores.


\textbf{Palabras clave:} Optimización binivel, MPECs, Condiciones necesarias de optimalidad, Punto estacionario.

\chapter*{Abstract}
The bilevel optimization problem is defined as minimizing a function on a set determined by the optimal points of a mathematical programming model. The optimization at the lower level depends on the decisions made at the upper level, thus creating an interdependent relationship between both levels.

To address this problem, the lower-level problem is replaced by the necessary optimality conditions and the resulting (relaxed) optimization problem with complementarity constraints is solved.

This work proposes a way to generate two-level problems whose relaxed problem has a stationary point belonging to one of the following classes: strongly stationary, M-stationary, or C-stationary, depending on the multipliers.

\textbf{Keywords:} Bi-level optimization, MPECs, Necessary optimality conditions, Stationary point.

	
	%
\chapter*{Agradecimientos}

El camino recorrido hasta este momento ha sido posible gracias al apoyo y dedicación de personas extraordinarias que han estado presentes en cada paso de mi formación académica. Deseo expresar mi más sincero agradecimiento a todos los profesores que han contribuido a mi desarrollo profesional durante mi trayectoria universitaria.

Mi más profunda gratitud a la profesora Dra. C. Gemayqzel, mi tutora, por su paciencia, dedicación y capacidad para transmitir conocimientos en el fascinante mundo de la optimización binivel. Su apoyo y comprensión han sido fundamentales para la culminación exitosa de este trabajo.

A mi familia, quiero agradecerles especialmente por su infinita paciencia y dedicación en los momentos más difíciles de este proceso. Su apoyo constante y su aliento para no abandonar nunca me han motivado a seguir adelante y alcanzar mis metas.

A mis amigos y compañeros de estudio, gracias por su tiempo, dedicación y palabras de aliento, que enriquecieron esta experiencia académica. 

Un agradecimiento especial a mis mascotas, por su compañía y alegría en los momentos desafiantes, y a ChatGPT, cuya asistencia fue invaluable en la redacción y organización de mis ideas.

Este trabajo no habría sido posible sin el soporte y la guía de cada uno de ustedes. Gracias por ser parte fundamental de este capítulo de mi vida.




	
	\tableofcontents
		
	\mainmatter
	\include{Introducción.tex}
		
	%\chapter{Preliminares}


La optimización binivel es un problema de optimización donde un subconjunto de variables debe ser la solución óptima de otro problema de optimización, parametrizado por las variables restantes. Este problema tiene dos niveles jerárquicos de decisión: el nivel superior (líder) y el nivel inferior (seguidor).
Este tiene dos características principales: primero, el problema del nivel inferior actúa como una restricción para el nivel superior; segundo, la solución del nivel inferior depende de las variables del nivel superior, creando una interdependencia entre ambos niveles. Por ello, el líder debe anticipar la respuesta óptima del seguidor al tomar decisiones.
%Explicación breve de binivel
En términos abstractos, la optimización binivel busca minimizar una función objetivo de nivel superior, $F(x, y)$, donde $x$ son las variables de decisión del líder y $y$ son las variables del seguidor. 

% Explicación del capitulo
En este capítulo se abordará los conocimientos necesarios para el desarrollo de esta tesis. Consta de secciones sobre los conceptos fundamentales de la Optimización Binivel del Caso Optimista, así como su reformulación KKT,
además de nociones clave sobre los Problemas Matemáticos con Restricciones de Equilibrio (MPEC), 
y por último la modelación de los problemas binivel en el lenguaje de programación Julia y métodos seleccionados implementados en este lenguaje para su resolución. 

\section{Optimización Binivel }
%Intro de la sección
El problema de optimización binivel en general se define de la siguiente manera
\begin{equation} \label{eq:Def1Binivel}
    \begin{array}{l}
       \min_{x} \quad F(x, y)\\
        s.a \left\{ \begin{array}{l}
            x \in T \\
             y \in S(x) = \arg  \min_{y} \{ f(x, y) \quad s.a \quad y \in  \cal{H}\\
            x,y \in M^0 
        \end{array}\right.
        \end{array} \end{equation}
En la mayoria de las aplicaciones, los problemas de optimizacion del nivel inferior y superior son modelos de programacion matematica, por lo  que tienen la estructura  *** con S(x)
\begin{equation}
\begin{aligned}
& \min_{x} \; F(x, y) \\
& \text{sujeto a } \\
& G(x) \leq 0, \\
& y \in \argmin_{y} \left\{ f(x, y) \mid v(x, y) \leq 0 \right\}
\end{aligned}
\label{eq:ProblemaGeneral}
\end{equation}
Para reducir las notaciones asumimos que en ambos problemas solo hay restricciones de desigualdad. O sea, se considerara el modelo 
\begin{equation}
\begin{aligned}
& \min_{x} \; F(x, y) \\
& \text{sujeto a } \\
& G(x) \leq 0, \\
& y \in \argmin_{y} \left\{ f(x, y) \mid v(x, y) \leq 0 \right\}
\end{aligned}
\label{eq:ProblemaGeneral}
\end{equation} 
%TODO:
% ***CON S(x)
%Para el concepto de solucion es necesario tener en cuenta la decision de ... parrafo de intro cambiado
%En este sentido, se consideran dos enfoques principales: el optimista y el pesimista.  El problema optimista consiste en  y se 
%% Formulación general de optimización binivel


% DIFERENCIA OPTIMISTA PESIMISTA
% Decir que en el libro de Dempe se hablan de estos dos enfoques
%Enfoque optimista
En el enfoque optimista, se asume que el seguidor, que actúa en el nivel inferior, elegirá la solución más favorable para el líder, quien toma decisiones en el nivel superior. Este es considerado más tratable y, en ciertas situaciones favorables, puede simplificarse a un problema convexo. Además, en el contexto de múltiples objetivos, el enfoque optimista permite alcanzar el mejor frente de Pareto posible, ver \cite{DempeyZemkoho2020}.

%\begin{center}   
%    Bajo el enfoque optimista :
%    \end{center}
    % Enfoque optimista (selección favorable)
    \begin{equation}
    \begin{aligned}
    & \min_{x} \; \inf_{y \in S(x)} F(x, y) \\
    & \text{con } S(x) = \argmin_{y} \left\{ f(x, y) \mid v(x, y) \leq 0 \right\}
    \end{aligned}
    \label{eq:EnfoqueOptimistaGeneral}
    %\caption*{Problema de optimización bajo el enfoque optimista}
    \end{equation}
    \captionof{figure}{Problema de optimización bajo el enfoque optimista.}

% Enfoque pesimista
Por otro lado, el enfoque pesimista asume que el seguidor seleccionará la opción menos favorable para el líder entre las soluciones óptimas disponibles, el cual es más complejo de resolver y puede incluso no tener solución. A menudo, se requieren reformulaciones para abordar estos problemas, lo que lo convierte en un reto teórico y computacional significativo en situaciones de múltiples objetivos, conduce al peor frente de Pareto posible, ver \cite{Sinha2017ARO}.
%\begin{center}
%    Bajo el enfoque pesimista:
%\end{center}
% Enfoque pesimista (peor caso)
\begin{equation}
\begin{aligned}
& \min_{x} \; \sup_{y \in S(x)} F(x, y) \\
& \text{con } S(x) = \argmin_{y} \left\{ f(x, y) \mid v(x, y) \leq 0 \right\}
\end{aligned}
\label{eq:EnfoquePesimistaGeneral}
\end{equation}
\captionof{figure}{Problema de optimización bajo el enfoque pesimista.}
%Decir que el optimista es mejor
Es relevante destacar que la mayoría de la literatura se centra en el enfoque optimista debido a su mayor facilidad de tratamiento. Sin embargo, el otro también tiene su utilidad, especialmente en la modelación de situaciones donde se considera la aversión al riesgo, ver \cite{DempeyZemkoho2020}. En este contexto, los términos ''líder'' y ''seguidor'' se utilizan para describir los roles en el modelo a optimizar; el líder toma decisiones considerando las posibles reacciones del seguidor, quien a su vez reacciona seleccionando su mejor opción, ver \cite{Sinha2017ARO}.

% Decir que se va ha hablar del enfoque optimista
Dado que en la tesis trataremos sobre problemas binivel de enfoque optimista mostraremos algunos resultados referidos al modelo resultante.


\newpage
% Definición de problema binivel

\textbf{Optimización Binivel Optimista.} Un problema de optimización binivel optimista, se formula como:

\begin{align}
    \min_{x , y} & \quad F(x, y) \notag \\
    \text{s.t.} & \quad g_i(x, y) \leq 0, i=1\ldots q,  \notag \\
    & \quad h_i(x,y) = 0, i=1 \ldots r, \notag\\
    & \quad y \in S(x), \notag\\
    &\intertext{donde $S(x)$ es el conjunto de soluciones óptimas del problema parametrizado por $x$} \tag{\theequation}\\
    \min_{y \in Y(x)} & \quad f(x, y) \notag \\
    \text{s.t.} & \quad v_j(x, y) \leq 0, j=1\ldots s, \notag\\
    & \quad u_j(x,y) = 0, j=1 \ldots t, \notag\\
    %&\text{Definición de un Problema Binivel Optimista} \notag\\
\end{align}
\label{eq:DefBinivelOptimista} 
%\textbf{Donde $M^0(x,y)$ es el conjunto de restricciones comunes para ambos niveles}

\begin{samepage}
%Dimensiones
donde 
% Dimension de x and y
\begin{equation*}
    x \in R^{n},\quad \quad y \in R^{m}
\end{equation*}
\begin{align*}
    %Dimension de F 
& F(x,y) : \mathbb{R}^{n} \times \mathbb{R}^{m} \to \mathbb{R},
%Dimension g_i
& \quad  g_i(x,y)  : \mathbb{R}^{n} \times \mathbb{R}^{m} \to \mathbb{R} ,
%Dimension h_i
& \quad h_i(x,y)  : \mathbb{R}^{n} \times \mathbb{R}^{m} \to \mathbb{R} ,
\\
% Def nivel inferior
%Dimension de f_i
&f(x,y) : \mathbb{R}^{n} \times \mathbb{R}^{m} \to \mathbb{R},
% Dimension de v_j
& \quad v_j(x,y)  : \mathbb{R}^{n} \times \mathbb{R}^{m} \to \mathbb{R} ,
% Dimension u_j
& \quad u_j(x,y)  : \mathbb{R}^{n} \times \mathbb{R}^{m} \to \mathbb{R} ,
\end{align*} 
\end{samepage}

%\begin{equation*}
%    %Dimension de M_i
%m_i \in M^0(x,y) \| m_i : f \mathbb{R}^{n} \times \mathbb{R}^{m} \to \mathbb{R}
%\end{equation*}

% Esto dejarlo en una sola pagina
%TODO: RESTRICCIONES DEL LIDER Y DEL SEGUIDOR NO IMPLICITA y EXPLICITA
 %TODO: RESTRICCIONES DEL LIDER Y DEL SEGUIDOR NO IMPLICITA y EXPLICITA
\newpage


Esta minimización está sujeta a dos tipos de restricciones: 
%Restricciones explicación old
%las restricciones explícitas para el líder, $x \in X$, donde $X$ es el conjunto de valores factibles para las variables del líder; y las restricciones implícitas impuestas por el seguidor, donde $y$ debe pertenecer al conjunto de soluciones óptimas del problema de optimización del seguidor, $\arg\min\{f(x, y) : y \in Y(x)\}$. En este contexto, $f(x, y)$ es la función objetivo del nivel inferior, y $Y(x)$ representa las restricciones del nivel inferior, las cuales pueden depender de las variables de decisión del líder, $x$.
% Restricciones nuevas
\begin{itemize}
    \item \textbf{Restricciones explícitas:}
    \begin{itemize}
        \item Desigualdades del líder: \( g_i(x,y) \leq 0 \quad \forall i \)
        \item Igualdades del líder: \( h_i(x,y) = 0 \quad \forall i \)
    \end{itemize}
    
    \item \textbf{Restricciones implícitas:}
    \begin{itemize}
        \item Impuestas por el seguidor: \( y \in \arg\min\{f(x, y) : v_j(x,y) \leq 0,\ u_j(x,y) = 0\} \),\\
        donde \( f(x, y) \) es la función objetivo del seguidor
    \end{itemize}
    
	 
	\item También se puede asumir que ambas variables tienen restricciones conjuntas, las restricciones $ g_i(x,y), \quad h_i(x,y), \quad v_j(x,y), \quad u_j(x,y) $ que dependen de las variables $x$ y $y$.
	
	
\end{itemize}
 % Hablar sobre el enfoque optimista
%En el caso del enfoque optimista si para el nivel inferior existen más de un punto que resuelve el problema del nivel inferior y tomará la que más beneficie al nivel superior. 
 % Decir que solo restricciones desigualdad
Por simplicidad solo se consideran restricciones de desigualdad.


%\subsection{Transformación de los problemas de dos niveles}
		
		Los problemas de dos niveles pueden ser reformulados en un problema de un solo nivel al reemplazar el problema del nivel inferior por las condiciones KKT de este en las restricciones del primer nivel. 
		
		%Para el caso de los SLSMG donde se tiene un problema de optimización en el nivel inferior este sustituye por de las condiciones KKT de este, obteniendo un MPEC \autocite{aussel2020}.
        
% Descripción del modelo en KKT
       \begin{table}[H]
        \centering
        \begin{equation}
            \begin{array}{l}
                \underset{\substack{x, y, \lambda_j}}{\min} \quad F(x, y)\\
                s.a \left\{ 
                \begin{array}{l}
                    g_i(x, y) \leq 0, i=1\ldots q,\\
                    \nabla_{y} f(x, y) + \sum_{j=1}^{s} \nabla_{y} v_j(x, y) \lambda_j = 0, \\
                    v_j(x, y) \leq 0, j=1\ldots s,\\
                    v_j(x, y)\lambda_j = 0, j=1\ldots s, \\
                    \lambda_j \geq 0, j=1\ldots s\\
                \end{array}\right.
                \tag{\theequation}
            \end{array}
            \label{eq:KKT_Optimista}
        \end{equation}
        \caption*{MPEC resultante}

    \end{table}
    
% Añadir aclaratoria
Los tres últimos grupos de restricciones expresan que $v$ y $\lambda$ están restringidas en signo y que al menos una es 0. Estas condiciones son conocidas como \textbf{restricciones de complementariedad}. Estos modelos corresponden a la clase de problemas de programación matemática con restricciones de complementariedad (MPEC). A continuación, presentamos los resultados de esta área necesarios para el desarrollo de esta tesis.
%TODO: Como se dijo en la intro existe otra forma, en un problema de programacion matematica  usando funcion del valor extremal sin embargo salvo en casos muy particulares esta no tiene una expresion explicita. La ventaja dsel kkt por mpec es que es una forma explicita, pero se requiere algun tipo de condicion de regularidad del conjunto como LICQ independencia  de v activa, que caso ded no cumpliser poder existir puntos optimos que no lo recoga el kkt. 

\section{Sobre los Programas Matemáticos con Restricciones de Equilibrio (MPEC)}
%Aca va la sección que explica los MPEC
% Que es un MPEC
Un \textbf{Programa Matemático con Restricciones de Equilibrio (MPEC)} es un tipo de programa no lineal que incluye restricciones de equilibrio, específicamente, restricciones de complementariedad.

%Formula MPEC Genérico
\begin{equation}
\begin{aligned}
\min \quad & f(z)  \\
\text{s.t.} \quad & g(z) \leq 0, \quad h(z) = 0 \\
& G(z) \geq 0, \quad H(z) \geq 0, \quad G(z)^T H(z) = 0 \\
&\text{Definición de MPEC} \\
\end{aligned}  
%\tag{\theequation} 
\label{eq:DefMpec}
\end{equation}

%Dimensiones MPEC Genérico
donde $f: \mathbb{R}^n \to \mathbb{R}$, $g: \mathbb{R}^n \to \mathbb{R}^m$, $h: \mathbb{R}^n \to \mathbb{R}^p$, $G: \mathbb{R}^n \to \mathbb{R}^\ell$, y $H: \mathbb{R}^n \to \mathbb{R}^\ell$ son funciones continuamente diferenciables. Debido al término de complementariedad en las restricciones, los programas de este tipo son algunas veces referidos como programas matemáticos

Estas restricciones de complementariedad se expresan como: 

%Ecuación para poner que son las restricciones de complementariedad
\begin{equation}
    G(z) \geq 0, \quad H(z) \geq 0, \quad \text{y} \quad G(z)^T H(z) = 0 \label{eq:RestriccionesComplementariedadAbstracto}
\end{equation}

Los MPEC incluyen restricciones donde el producto de dos funciones (\textit{G} y \textit{H}) debe ser cero, y ambas funciones deben ser no negativas. Esto se conoce como restricción de complementariedad. Debido a las restricciones de complementariedad, los MPEC no cumplen con las condiciones de regularidad estándar, lo que hace que las condiciones de KKT no sean directamente aplicables como condiciones de optimalidad de primer orden. Los MPEC son utilizados para modelar problemas donde existen restricciones de equilibrio, como problemas de ingeniería y economía, ver \cite{Flegel2003AFJ,DempeyZemkoho2020}.


%\subsection{Resultados sobre MPECs}

%Notación de los indices activos del documento
A continuación como en \cite{Flegel2003AFJ}, se introduce las siguientes definiciones.
\begin{definition}[Conjunto de Índices]
Dado un vector factible $z^*$ del MPEC \eqref{eq:DefMpec}, definimos los siguientes \textit{conjuntos de índices}:
\begin{equation}
\begin{aligned}
J_G &:= J_G(z^*) := \{i|G_i(z^*) = 0, H_i(z^*) > 0\} \\
J_{GH}:= J_{GH}:(z^*) := \{i|G_i(z^*) = 0, H_i(z^*) = 0\}  \\
J_H &:= J_H(z^*) := \{i|G_i(z^*) > 0, H_i(z^*) = 0\}  \\
\end{aligned}
\label{eq:ConjuntoDeIndices} 
\end{equation}
\end{definition}

% Definición de TNLP Problema No Lineal Ajustado
Para definir calificaciones de restricción alteradas, introducimos el siguiente problema, dependiente de $z^*$:

\begin{definition}[Programa No Lineal Ajustado (TNLP)]

Un \textit{Programa No Lineal Ajustado $TNLP := TNLP(z^*)$} es:
%TODO: INcomporar indices y rango en que se mueven
\begin{equation}
\begin{aligned} 
\min \quad & f(z) \\
\text{s.t.} \quad & g_{i}(z) \leq 0 \quad h(z) = 0  \\
& G_{\alpha \cup \beta}(z) = 0 \quad G_{\gamma}(z) \geq 0  \\
& H_{\alpha}(z) \geq 0 \quad H_{\gamma \cup \beta}(z) = 0  \\
& \text{Problema No Lineal Ajustado (TNLP)}  
\end{aligned}
\label{eq:ProblemaAbstractoTNLP}
\end{equation}
    
\end{definition}

El TNLP \eqref{eq:ProblemaAbstractoTNLP} puede ahora ser usado para definir variantes MPEC adecuadas de las calificaciones de restricción estándar de independencia lineal, Mangasarian-Fromovitz y Mangasarian-Fromovitz estricta (LICQ, MFCQ, y SMFCQ en su forma abreviada).

\begin{definition}[MPEC-LICQ]
El MPEC \eqref{eq:DefMpec} se dice que satisface la \textit{MPEC-LICQ} (MPEC-MFCQ, MPEC-SMFCQ) en un vector factible $z^*$ si el correspondiente TNLP$(z^*)$ satisface la LICQ (MFCQ, SMFCQ) en ese vector $z^*$.
\end{definition}
    
% Hacer notar que un mínimo de z^* implica la existencia de \lambda^* ....
%En este punto, es importante notar que bajo MPEC-MFCQ, un mínimo local $z^*$ del MPEC \eqref{eq:DefMpec} implica la existencia de un multiplicador de Lagrange $\lambda^*$ tal que $(z^*, \lambda^*)$ satisface las condiciones KKT para el programa \eqref{eq:ProblemaAbstractoTNLP}. Por lo tanto, si asumimos que MPEC-MFCQ se cumple para un mínimo local $z^*$ del MPEC \eqref{eq:DefMpec}, podemos usar cualquier multiplicador de Lagrange $\lambda^*$ (que ahora se sabe que existe) para definir el MPEC-SMFCQ, es decir, tomando $(z^*, \lambda^*)$.

%Definición de los puntos estacionarios del MPEC
% Flegel and Kanzow 2003
En el contexto de los MPEC en \cite{Flegel2003AFJ} se exponen varios tipos de puntos estacionarios que son cruciales para analizar la optimalidad, los cuales son los siguientes:
\begin{definition}[Punto Factible]
    Un \textit{punto factible} $z$ del MPEC se llama débilmente estacionario si existe un multiplicador de Lagrange $\lambda = (\lambda^g, \lambda^h, \lambda^G, \lambda^H)$ tal que se cumplen las siguientes condiciones:
    
%TODO: USar multiplicadores analogos a como esta en el bijnivel de miu lambda beta 
% EN este orde  mui-g, alpha-h beta-G gamma-H
\begin{align}
& \nabla f(z) + \sum_{i=1}^m \lambda_i^g\nabla g_i(z) + \sum_{i=1}^p \lambda_i^h\nabla h_i(z) - \sum_{i=1}^{\ell} [\lambda_i^G\nabla G_i(z) + \lambda_i^H\nabla H_i(z)] = \vec{0} \notag\\
    & \quad \begin{aligned}
        & \lambda_\alpha^G \text{ libre}, \quad \lambda_\beta^G \text{ libre}, \quad \lambda_\gamma^G = 0, \\
        & \lambda_\gamma^H \text{ libre}, \quad \lambda_\beta^H \text{ libre}, \quad \lambda_\alpha^H = 0, \\
        & g(z) \leq 0, \quad \lambda^g \geq 0, \quad (\lambda^g)^T g(z) = 0.
    \end{aligned} \\
& \label{eq:Definicion_punto_debilemente_estacionario} \notag
\end{align}
\end{definition}

% Definición de puntos estacionarios
Este concepto de estacionariedad, sin embargo, es una condición relativamente débil. Existen conceptos más fuertes de estacionariedad que se derivan y estudian en otros lugares. En particular, se tienen las siguientes definiciones:
% un punto factible $z$ con el correspondiente multiplicador de Lagrange $\lambda = (\lambda^g, \lambda^h, \lambda^G, \lambda^H)$ se llama:

\begin{itemize}
% C-estacionario
\item \begin{definition}[Punto C-estacionario]
  El punto $z$ es: \textit{C-estacionario} si, para cada $i \in \beta$, $\lambda_i^G\lambda_i^H \geq 0$ se cumple.
\end{definition}
%M-Estacionario
\item \begin{definition}[Punto M-estacionario]
    El punto $z$ es: \textit{M-estacionario} si, para cada $i \in \beta$, o bien $\lambda_i^G,\lambda_i^H > 0$ $\vee$ $\lambda_i^G \lambda_i^H = 0$.
\end{definition}
%Fuertemente estacionario
\item \begin{definition}[Punto Fuertemente estacionario]
    El punto $z$ es: \textit{fuertemente estacionario} o \textit{estacionario primal-dual} si, para cada $i \in \beta$, $\lambda_i^G, \lambda_i^H \geq 0$.
\end{definition}
% A-estacionario
%\item \begin{definition}[Punto A-estacionario]
%   Sea $z^*$ un punto débilmente estacionario del MPEC \eqref{eq:DefMpec} será \textit{A-estacionario} si existe un multiplicador de Lagrange correspondiente $\lambda^*$ tal que:
%\begin{equation}
%(\lambda_i^G)^* \geq 0 \quad \vee \quad (\lambda_i^H)^* \geq 0 \quad \forall i \in \beta \notag
%\end{equation}
%\end{definition}
\end{itemize}

%TODO:
Estas definiciones son condiciones necesarias de optimalidad bajo ciertas condiciones de regularidad, las m;as fuerte basda en la MPEC-LICQ conlleva al siguiente resultado:
% Teorema que si z^* es min local del MPEC Si se cumple la LICQ entonces existe .....
\begin{theorem} 
Sea $z^* \in \mathbb{R}^n$ un mínimo local del MPEC \eqref{eq:DefMpec}. Si MPEC-LICQ se cumple en $z^*$, entonces existe un único multiplicador de Lagrange $\lambda^*$ tal que $(z^*, \lambda^*)$ es \textit{fuertemente estacionario}.
\end{theorem}

% Modelación en Julia
% Se explica de que va BilevelJump
\section{Modelación en Julia}
% Que es BilevelJump
BilevelJuMP.jl es un paquete de Julia diseñado para modelar y resolver problemas de \textbf{optimización binivel}, también conocidos como problemas de optimización de dos niveles o jerárquica, para más detalles ver \cite{BilevelJump}.
Estos problemas se caracterizan por tener dos niveles de decisión: un nivel superior y un nivel inferior, donde las decisiones del nivel superior influyen en las decisiones del nivel inferior, y viceversa \cite{BilevelJump}.
% Que problemas resuelve
Este paquete permite abordar una amplia variedad de tipos de problemas. En este trabajo usaremos las siguientes bibliotecas:
% Añadir explicación MPEC
\subsection*{JuMP}
% TODO: PONEWr esto  Permite gestionar una variedad de restricciones compatibles con JuMP, tales como restricciones cuadráticas, no lineales y enteras.
%Facilita la transformación de problemas binivel en problemas de programación matemática con restricciones de equilibrio (MPEC) y ofrece métodos para abordar las restricciones de complementariedad que surgen en estos casos.
  %TOD: Poner esto despues de expluicar las reformulaciones
\subsection*{BilevelJuMP}
BilevelJuMP.jl facilita la modelación de problemas que pueden representarse en la sintaxis de JuMP, ver documentación en \cite{JuMPPaper}, incluyendo restricciones lineales y no lineales, variables continuas y enteras, y diferentes tipos de objetivos.
\begin{itemize}
    \item \textbf{Problemas de optimización binivel generales:} BilevelJuMP.jl facilita la modelación de problemas que pueden representarse en la sintaxis de JuMP, ver documentación en \cite{JuMPPaper}, incluyendo restricciones lineales y no lineales, variables continuas y enteras, y diferentes tipos de objetivos.
    
    \item \textbf{Problemas con restricciones de determinados tipos en el nivel superior:} Permite gestionar una variedad de restricciones compatibles con JuMP, tales como restricciones cuadráticas, no lineales y enteras.
    
%TODO: MOdificar esto para decir que hace esto.
    
    
    \item \textbf{Problemas con diferentes tipos de reformulaciones:} Los usuarios pueden experimentar con diversas reformulaciones para las restricciones de complementariedad en los problemas MPEC, incluyendo SOS1, restricciones de indicador, Fortuny-Amat y McCarl (Big-M), entre otros.
    %TODO: POner que es comun a los dos 
    \item \textbf{Problemas que requieren solvers MIP y NLP:} BilevelJuMP.jl puede utilizar tanto solucionadores de \textbf{programación lineal mixta entera (MIP)} como solucionadores de \textbf{programación no lineal (NLP)}, dependiendo de las características del problema y la reformulación elegida.
\end{itemize}

%\subsection{Limitaciones del BilevelJuMP.jl}

A pesar de sus capacidades, BilevelJuMP.jl presenta algunas limitaciones, entre las cuales se encuentran:

\begin{itemize}
    \item Enfrentar dificultades en problemas altamente no lineales o con estructuras de optimización complejas que no se puedan representar adecuadamente en la sintaxis de JuMP.
    \item Existencia de ciertas restricciones que podrían no ser compatibles o que requieren transformaciones adicionales que podrían complicar el modelo.
    \item El problema es de gran escala afectando el rendimiento del solver, el cual puede ser un factor crítico.
    \item La formulación y resolución de problemas muy específicos o especializados podrían no estar completamente optimizadas en el paquete.
\end{itemize}


% Explicación algoritmos 
\section{Métodos de Reformulación para Optimización Binivel}
La optimización binivel presenta desafíos particulares debido a su naturaleza jerárquica y las condiciones de complementariedad resultantes. A continuación, se presentan los principales métodos de reformulación implementados en la literatura, basados en la transformación del MPEC \eqref{eq:KKT_Optimista}.
\subsection{Método Big-M}

El método Big-M (Fortuny-Amat y McCarl)  es una técnica fundamental para reformular problemas de optimización binivel en problemas MPEC. 
Este método aborda específicamente las condiciones de complementariedad que surgen en estas reformulaciones, transformando el problema original en un problema de programación lineal mixta entera (MILP).

La reformulación mediante Big-M introduce un parámetro M suficientemente grande y variables binarias para transformar las condiciones de complementariedad no lineales en restricciones lineales. Para una condición de complementariedad 
$v_j(x,y)\lambda_j = 0$ en \eqref{eq:KKT_Optimista}, el método introduce una variable binaria  $f_j \in \{0,1\}$ y cotas superiores $M_p$, $M_d$ bajo  las siguientes restricciones:

\begin{align}
    v_j(x,y) &\geq -M_p(1 - f_j) \notag \\
    \lambda_j &\leq M_d f_j \\
	%TODO: EN vez sde f_j poner \delta_j
    f_j &\in \{0,1\} \notag
\end{align}

donde $M_p$ y $M_d$ son valores grandes para las variables primales y duales, respectivamente. La efectividad del método depende crucialmente de la selección apropiada de estos valores, que deben ser suficientemente grandes para no excluir la solución óptima, pero no excesivamente grandes para evitar inestabilidades numéricas \cite{BilevelJump}.


\subsection{Método SOS1}

El método de Conjuntos Ordenados Especiales tipo 1 (SOS1) evita el uso de parámetros Big-M mediante restricciones de tipo conjunto. Para cada par complementario $(v_j, \lambda_j)$:

\begin{equation}
%DONDe SOS1 es v_j * \lambda_j 
[v_j(x,y); \lambda_j] \in \text{SOS1} \label{eq:SOS1_reform}
\end{equation}

Esta restricción fuerza que al menos una variable en el par sea cero, preservando la no linealidad original sin necesidad de cotas. Es particularmente eficaz en problemas con estructura cónica \cite{BilevelJump}.


\subsection{Método ProductMode}

El método ProductMode representa un enfoque directo para manejar las condiciones de complementariedad en su forma de producto original. Este método es particularmente útil cuando se trabaja con solucionadores de programación no lineal (NLP), aunque no garantiza la optimalidad global.

La implementación del ProductMode mantiene la restricción de complementariedad en su forma original:

\begin{equation}
    v_j(x,y) \cdot \lambda_j \leq t \label{eq:ProductMode_reg}
\end{equation}

donde $t > 0$ es un parámetro de regularización pequeño. Esta formulación, aunque no satisface las condiciones de calificación de restricciones estándar, es útil para obtener soluciones iniciales y puede ser especialmente efectiva cuando se combina con solucionadores NLP, ver \cite{BilevelJump}.

% TODO: hjacer suempre cierre capitulo, y dicinedo que se dijo y motivar

%\subsection{Comparación de Métodos}
%La Tabla \ref{tab:comparacion_metodos} resume los requisitos y características de cada método:
%
%\begin{table}[H]
%\centering
%\begin{tabular}{l|l|l}
%\textbf{Método} & \textbf{Solver Requerido} & \textbf{Ventajas} \\ \hline
%Big-M & MIP & Estabilidad numérica controlada \\
%SOS1 & MIP con SOS1 & Sin parámetros ad-hoc \\
%ProductMode & NLP & Manejo de no linealidades \\
%\end{tabular}
%\caption{Comparación de métodos de reformulación}
%\label{tab:comparacion_metodos}
%\end{table}
%
%Para casos con restricciones no lineales $v_j(x,y)$, se recomienda combinar ProductMode con técnicas de linealización por tramos \cite[Apéndice B]{BilevelJump}. Todos estos métodos están implementados en BilevelJuMP.jl, permitiendo experimentación con múltiples reformulaciones \cite[Sección 4]{BilevelJump}.
%
%\section{Add}
%\begin{equation}
%    \min_{x,y,\lambda_j,f_j} F(x,y)
%\end{equation}
%
%sujeto a:
%\begin{align}
%    & g_i(x,y) \leq 0, && i = 1,\ldots,q, \\
%    & \nabla_y f(x,y) + \sum_{j=1}^s \nabla_y v_j(x,y) \lambda_j = 0, \\
%    & v_j(x,y) \leq 0, && j = 1,\ldots,s, \\
%    & \lambda_j \geq 0, && j = 1,\ldots,s, \\
%    & v_j(x,y) \geq -M_p \cdot (1-f_j), && \text{(no lineal si } v_j \text{ es no lineal)} \\
%    & \lambda_j \leq M_d \cdot f_j, && j = 1,\ldots,s, \\
%    & f_j \in \{0,1\}, && j = 1,\ldots,s.
%\end{align}

	%\chapter{IMPLEMENTACIÓN}

En este capítulo se mostrará la forma y los pasos para generar el problema deseado en el contexto de un generador de puntos estacionarios para problemas binivel. Este capítulo aborda el diseño y la implementación del generador, explicando detalladamente las metodologías empleadas, los criterios necesarios para asegurar la validez de los puntos obtenidos y los algoritmos utilizados en cada etapa del proceso.

%\section*{Nomenclatura a utilizar:}
\begin{table}[H]
    \centering
    \caption{Nomenclatura a utilizar}
    \begin{tabular}{l m{360pt}}
        $ F(x,y) $              & Función del líder                                                                                                          \\
        $ g_i(x,y) $              & Restricciones de desigualdad del líder                                                                                                            \\
        $ h_i(x,y) $                 & Restricciones de igualdad del lider                       \\
        $ f(x,y) $           & Funcion del follower                                                               \\
        $ v_j(x,y) $              &  Restricciones de desigualdad del follower  \\
        $ u_j(x,y) $     & restricciones de igualdad del follower\\
        $ v_{j}^{\star} $    & Restriccion de disigualdad del follower activa despues de modificarse el problema            \\
        $z_{est}$         & Punto que se desea que sea estacionario\\
        $J_0^g$   & Índices activos del nivel superior \\
        $J_0^v$   & Índices activos del nivel inferior \\
        $ \alpha  $             & Vector de entrada de dimensión igual cantidad variables follower                                                                                                      \\
        $ \beta $               & Multiplicador asociado al $v_{j}^{\star}$ en el nivel superior          \\
        $ \lambda $              & Multiplicador asociado al $v_{j}^{\star}$ en el nivel inferior\\
        $\gamma$                & Valor asociado a cada $v_{j}^{\star}$ con $\alpha=\vec{0}$\\
    \end{tabular}

    \end{table}


\section{Generación del problema}
A continuacion se mostrarán los pasos a seguir para generar el problema deseado donde el punto sea estacionario de la clase específica.
\subsection{Requerimientos de entrada}

Para generar un problema es necesario inicialmente introducir un problema de optimización binivel del tipo SLSF, un punto ($z_{est}$) el cual
se desea que sea estacionario y los multiplicadores con signo o valores con ciertas condiciones en dependencia del tipo de estacionariedad requerida.

El problema de entrada tiene que cumplir con ser un programa de optimizacion binivel como \refeq{eq:DefBinivelOptimista}. 
Debe de conocerse sus índices activos de las restricciones de desigualdad en cada nivel:
\begin{itemize}
    % Leader
    \item Nivel Superior:
            \begin{equation}
             J_0^G=\{i | g_i(x,y)=0\}
            \label{J_0_level_superior} %al conjunto restricciones de desigualdad que sean índices activos del nivel superior
            \end{equation}\\
    Cuyos índices activos tienen un $\mu_i$ relacionados.
        
    %Follower
    \item Nivel Inferior:
                Se tiene en cuenta con respecto a los valores de $\lambda_i$ y $v_j(x,y)$.
                
                \begin{itemize}
                    %J_0_v_l_0
                    \item \begin{equation}
                        J_1^v=\{j | v_j(x,y)=0 \land \lambda=0 \} %al conjunto restricciones de desigualdad que sean índices activos
                        \label{J_0_lambda_0_level_inferior}
                        \end{equation} 
                    %J_0_v_l_positiva
                    \item \begin{equation}
                        J_2^v=\{j | v_j(x,y)=0 \land \lambda>0 \}
                        \label{J_0_lambda_pos_level_inferior}
                    \end{equation}
                    %J_Ne_L0_v
                    \item \begin{equation}
                        J_3^v=\{j | v_j(x,y)< 0 \land \lambda=0 \}
                        \label{J_neg_lambda_0_level_inferior}
                    \end{equation}
                \end{itemize}
    Cada índice activo tiene multiplicadores $\beta_j$ y en dependencia del caso $\gamma_j$ asociados.
\end{itemize}

% Condiciones
%Hablar sobre conocerr vector \alpha
Se debe conocer el valor del $\vec{\alpha}$. 
Además para cada $j_i \in J_0^G$ \refeq{J_0_level_superior} debe de asignarse su $\mu_i \geq 0$ correspondiente,
para el nivel inferior debe tenerse en cuenta si $\alpha=\vec{0}$ en cuyo caso afirmativo $\gamma_j$ es un valor de entrada

% COndiciones \Beta_i \Gamma_i para ptos estacionarios
Para obtener las diferentes clases de puntos estacionarios debe de introducirse $\beta_j$, $\gamma_j$, en caso de ser este último valor de entrada tal que, en los índices activos $J_3^v$ \refeq{J_0_lambda_0_level_inferior}: 
\begin{itemize}
    %C estacionario
    %Betai, gammai>0, betai,gammai<0 o 
    %Betai libre gammai=0
    %Gammai libre , betai=0
    \item \textbf{Punto C-Estacionario:}\\
    \begin{equation}
        \beta_j * \gamma_i >=0
        \label{Requisitos puntos C-Estacionario}
    \end{equation}
    %M estacionario
    %Betai, gammai>0,
    %Betai libre gammai=0
    %Gammai libre , betai=0
    \item \textbf{Punto M-Estacionario}\\
    \begin{equation}
        \begin{aligned}
            &\beta_j \land \gamma_j>0\\
            &\beta_j \quad \text{Libre} \quad \gamma_j=0\\
            &\beta_j=0 \quad \gamma_j \quad \text{libre}\\
        \end{aligned}
        \label{Requisitos puntos M-Estacionarios}
    \end{equation}
    %Fuertemente estacionario
    %Betai, gammai>0,
    \item \textbf{Punto Fuertemente Estacionario}\\
     \begin{equation}
        \beta_j \land \gamma_j >0
    \label{Requisitos puntos Fuertemente Estacionarios}
    \end{equation}    
\end{itemize}
% Especificar que en caso de  que alpha no se el vector nulo se asumme que gamma_i es 0
\textbf{En caso que $\vec{\alpha} \neq \vec{0}$ se asume que $\gamma_j=0$ en las condiciones anteriores.}

% Explicar que se hace primero el punto factible
% Subseccion donde se explica que el punto se hace factible inicialmente

\subsection{Modificación del Problema Original}
Después de tener los datos de entrada necesarios se procede a la Modificación
del problema de entrada para que este sea estacionario del tipo requerido, $z_{est}$.

% Decir que en caso de que \alpha != de nulo entonoces se calcula b_j
Primero en caso de que $\alpha \neq 0$ en los $v_j(x,y) \in J_0^v$ se procede a calcular el $\vec{b_j}$ de la siguiente forma:

%Como calcular b_j
\begin{itemize}
    \item \textbf{ $v_j(x,y) \in J_2^v$ \refeq{J_0_lambda_pos_level_inferior}}:
        \begin{equation}
            b_j=  \frac{-\nabla_{y}{v_j(z_{est})}^T \cdot \alpha}{\|\mathbf{\alpha} \|_2}
        \end{equation}
    \item \textbf{$v_j(x,y) \in J_1^v \lor J_3^v$ (\refeq{J_0_lambda_0_level_inferior},\refeq{J_neg_lambda_0_level_inferior})}\\
    \begin{equation}
        b_j=  \frac{(-\nabla_{y}{v_j(z_{est})}^T \cdot \alpha)+\gamma_j}{\|\mathbf{\alpha} \|_2^2}
    \end{equation}
\end{itemize}


Después de realizado el calculo del $b_j$ se procede a modificar su indice activo correspondiente:

\begin{equation}
	v_{j}^{\star}(z_{est})=v_{j}(z_{est})+ ({b_j\alpha}^T)\cdot (y_1,y_2,\dots,y_m)
\end{equation}


% Explicar que se hace el punto factible en las restricciones
Luego en todas las restricciones del nivel inferior y superior
se evalua el punto y se verifica si es factible. En caso de no ser factible 
bajo las nuevas restricciones se añaden constantes $c_i$ en el caso de las restricciones del
nivel superior y $c_j$ para las del inferior. Esto conlleva a que sea un punto factible en el conjunto
de restricciones de ambos niveles sin perjudicar las propiedades relacionadas con la convexidad al ser una adicción
de funciones lineales. 

% Explicar que se realiza el KKT del nivel inferior
Al hacer factible se procede a realizar el KKT del nivel inferior con las modificiaciones anteriores evaluado en $z_{est}$.

Se plantea las condiciones KKT y se halla $\vec{bf}$ :

\begin{equation}
    \begin{aligned}
        &\nabla_{y}f(x,y)+\sum_{j=0}^{|V \in J_0|}(\lambda_j\nabla_{y}v_{j}^{\star}(x,y))+\vec{bf}=\vec{0}\\
        &\text{KKT del problema del nivel inferior}
    \end{aligned}
    \label{KKT_nivel_inferior}
\end{equation}
\linebreak
%Explicar que como se habia hallado ya el KKT del level inf ahora se puede realizar completo
Finalmente podemos construir la reformulación en MPEC como \refeq{eq:KKT_Optimista} calculando asi el $\vec{BF}$ y 
obteniendo el problema de forma reformulada.
% MPEC all KKT 
\begin{table}[H]
    \begin{equation}
        \label{KKT_del_MPEC}
    \end{equation}
	$\nabla_{xy}F(x,y)+\sum_{i=1}^{|G \in J_o|}(\mu_i\nabla_{xy}g(x,y))+[\nabla_{x,y}\nabla_{y}f(x,y)+\sum_{j=1}^{|V \in J_0|}\lambda_j\nabla_{xy}\nabla_{y}v_{j}^{\star}(x,y)]\alpha+\sum_{j=1}^{|V \in J_0|}(\beta_j\nabla_{xy}v_{j}^{\star}(x,y))+\vec{BF}=\vec{0}$  \label{KKT_del_MPEC}
\caption*{KKT del MPEC}
\end{table}

% Explicación algoritmica y manual de usuario

\section{implementación Algorítmica}

En la implementación computacional del generador propuesto en la sección anteriores
se utilizó el lenguaje de programación \textbf{Julia} \cite{Juliadocs} dado a su versatilidad
en expresiones y funciones matemáticas implementadas en su paquete base y sus bibliotecas externas como 
\textbf{BilevelJuMP} \cite{BilevelJump} que permite resolver problemas binivel SLSF lineales y cuadráticos
sin tener que transformar el problema original y \textbf{JuMP} \cite{JuMPPaper} el cual es una interfaz robusta para la optimización general,
todo ello brindando unas excelentes prestacioens de cómputo. Por ello se brindará una guía de usuario para el uso del generador.

\subsection{Guía de Usuario}
% Después compilar a ProblemGenerator
La implementación algorítmica se ha llevado a cabo mediante la creación de una biblioteca de Julia
llamada \textbf{ProblemGenerator}, con una sintaxis \textit{JuMP-like} intuitiva.
Primero debe tenerse un problema de optimización Binivel SLSF planteado como el \refeq{eq:DefBinivelOptimista}, el punto que se desea a ser estacionario ($z_{est}$) ,los indices activos descritos anteriormente 
según el caso así como sus multiplicadores correspondientes al tipo de estacionariedad requerida y el $\vec{\alpha}$.

%Dependencias para la biblioteca
\subsubsection{Dependencias}
%Se debe tener en cuenta que este paquete instalará las siguientes dependencias
\begin{itemize}
    \item Symbolics 
    \item LinearAlgebra
\end{itemize}

% Importar la Biblioteca
\subsubsection{Importación de la biblioteca}

\begin{lstlisting}[caption=Importar el Modulo]
    using ProblemGenerator
\end{lstlisting}

% Crear modelo base
\subsubsection{Crear el modelo base}
Debe llamarse a la función \textbf{GeneratorModel} para crear el modelo 
inicial 

% Explicar las funciones Upper y Lower
\subsubsection{Función para diferenciar los niveles}
En caso que se requiera como entrada un \textbf{Problem} se debe declarar
de que nivel se habla mediante evaluar el model en las funciones \textbf{Upper} para el nivel 
superior y \textbf{Lower} para el inferior.


%Crear para alpha nulo
Para crear un modelo donde $\vec{\alpha}=\vec{0}$
\begin{lstlisting}
    # Crear el modelo base del generador alpha=0
    model=GeneratorModel()
\end{lstlisting}

%Crear para alpha no nulo
Para crear un modelo donde $\vec{\alpha}\neq \vec{0}$
Se tiene que pasar como parámetro el valor del $\alpha$
el cual será un vector de Number.
\begin{lstlisting}
    # Crear el modelo base del generador alpha!=0
    alpha_vec=Vector::{Number}
    model=GeneratorModel(alpha_vec)
\end{lstlisting}

\subsubsection{Declarar Variables}
Las variables deben de 
declararse bajo la siguiente macro \textbf{@myvariables}, la cual 
recibe una funcion \textbf{Upper} o \textbf{Lower} que recibe el modelo
para las variables del nivel superior y las del nivel inferior respectivamente.

%Ejemplo de introducir al level superior
Ejemplo para introducir las variables del nivel superior
\begin{lstlisting}
    # Se declara en el nivel superior las variables x_1, x_2
    @myvariables Upper(model) x_1, x_2
\end{lstlisting}

%Ejemplo para introducir las variables del nivel inferior
Ejemplo para introducir las variables del nivel inferior
\begin{lstlisting}
   # Se declara en el nivel inferior las variables y_1, y_2
   @myvariables Lower(model) y_1,y_2
\end{lstlisting}


%Declarar las funciones objetivo
\subsubsection{Declarar Funciones Objetivo}

Para declarar las funciones objetivos debe asumirse que en ambos casos es un 
problema de minimización. Se utilizará la función \textbf{SetObjectiveFunction}
que recibe un \textbf{Problem} y una expresión de \textbf{Num}.

% Declaración de Función Objetivo de nivel Superior
Ejemplo dclaración una función objetivo del nivel superior.
Con: $$\min (x_1^2*y_1^2*y_2) + x_2$$
\begin{lstlisting}
    # Declarar la funcion objetivo del nivel superior
    # Min de ((x_1^2)*(y_1^2)*(y_2))+x_2
    SetObjectiveFunction(Upper(model),((x_1^2)*(y_1^2)*(y_2))+x_2)
\end{lstlisting}

%Declaración de Funcion objetivo Nivel inferior
Ejemplo dclaración una función objetivo del nivel inferior.
Con: $$\min (x_2^2*y_1^2*y_2)+x_1$$
\begin{lstlisting}
    # Declarar la funcion objetivo del nivel inferior
    # Min de ((x_2^2)*(y_1^2)*(y_2))+x_1
    SetObjectiveFunction(Lower(model),((x_2^2)*(y_1^2)*(y_2))+x_1)
\end{lstlisting}

% Definición de tipo de índices activos.
\subsubsection{Tipos de indices activos}
Antes de ilustrar como introducir las restricciones se va a explicar 
que los tipos de índices activos.

\begin{itemize}
    \item \textbf{Ambos Niveles}:
        \begin{itemize}
            \item \textbf{Normal} si no es un indice activo.
        \end{itemize}
    \item \textbf{Nivel Superior}:
     \begin{itemize}
        \item \textbf{J\_0\_g} Si es un índice como en \refeq{J_0_level_superior}.
     \end{itemize}
    \item \textbf{Nivel Inferior}:
    \begin{itemize}
        \item \textbf{J\_0\_LP\_v} Si es un índice como en \refeq{J_0_lambda_pos_level_inferior}.
        \item  \textbf{J\_0\_L0\_v} Si es un índice como en \refeq{J_0_lambda_0_level_inferior}.
        \item  \textbf{J\_Ne\_L0\_v} Si es un índice como en \refeq{J_neg_lambda_0_level_inferior}.
    \end{itemize}
\end{itemize}

%Ejemplo de como se debe de expresar en Julia
Los \textbf{RestrictionSetType} deben expresar en codigo de esta forma:
\begin{lstlisting}
    Normal 
    J_0_g 
    J_0_LP_v
    J_0_L0_v
    J_Ne_L0_v
\end{lstlisting}

Los \textbf{RestrictionSetType} so un enum de Julia.

% Declaracion de restricciones
\subsubsection{Declarar Restricciones}

Para declarar las restricciones se brindan dos funciones:
\begin{itemize}
    \item \textbf{SetLeaderRestriction:} Para el nivel superior.
    \item \textbf{SetFollowerRestriction:} Para el nivel inferior.
\end{itemize}
\begin{itemize}
    % Declarar sobre el nivel superior
    \item \textbf{Nivel Superior:}\\

% Declaracion para el nivel Superior
Para declarar las restricciones del nivel superior debe por cada restricción
llamarse a la función \textbf{SetLeaderRestriction} con:
\begin{itemize}
    \item El modelo (\textbf{model})
    \item La expresión de la restricción, del tipo \textbf{Num}
    \item El tipo de restriccion, del tipo \textbf{RestrictionSetType}
    \item El valor de $\mu_i$ correspondiente, del tipo \textbf{Number} 
\end{itemize}

Ejemplo para introducir la restricción:
\begin{align*}
    &x_1+y_2-y_1 \lq -9 \\
    &\text{Índice activo del tipo J\_0\_g \refeq{J_0_level_superior}}\\
    & \mu_i=0.3
\end{align*}

\begin{lstlisting}
    # Ejemplo de restriccion del nivel superior
    SetLeaderRestriction(model,x_1+y_2-y_1>9,J_0_g,0.3)
\end{lstlisting}

% Declarar sobre el nivel inferior
\item \textbf{Nivel Inferior:}\\
 Para declarar las restricciones del nivel inferior debe por cada restricción
llamarse a la función \textbf{SetFollowerRestriction} con:
\begin{itemize}
    \item El modelo (\textbf{model})
    \item La expresión de la restricción, del tipo \textbf{Num}
    \item El tipo de restriccion, del tipo \textbf{RestrictionSetType}
    \item El valor de $\beta_j$ correspondiente, del tipo \textbf{Number}
    \item El valor de $\lambda_j$ correspondiente, del tipo \textbf{Number}
    \item El valor de $\gamma_j$ correspondiente en caso de ser valor de entrada, del tipo \textbf{Number}
\end{itemize}

Ejemplo para introducir la restricción:

\begin{itemize}
    \item \textbf{$\vec{\alpha}\neq \vec{0}$:}\\
    \begin{align*}
        &((x_1^2)*(y_1^2))+x_2=0\\
        &\text{Índice activo del tipo J\_Ne\_L0\_v} \refeq{J_neg_lambda_0_level_inferior}\\
        & \beta_j=0.1\\
        & \lambda_j=0 \\
    \end{align*}
    \begin{lstlisting}
        # Para caso alpha!=0
        SetFollowerRestriction(model,((x_1^2)*(y_1^2))+x_2==0,J_Ne_L0_v,0.1,0)
    \end{lstlisting}
    \item \text{$\vec{\alpha}= \vec{0}$:}\\
    
    Análogo al caso anterior pero con $\gamma_j$ valor de entrada
    \begin{align*}
        &\gamma_j=0.4
    \end{align*}
    \begin{lstlisting}
        SetFollowerRestriction(model,((x_1^2)*(y_1^2))+x_2==0,J_Ne_L0_v,0.1,0,0.4)
    \end{lstlisting}
\end{itemize}

% Finalizar el itemize que tiene de explicacion como instroducir las restricciones de ambos niveles
\end{itemize}


\subsubsection{Introducir el Punto}
Ahora debe de introducirse el valor del punto que debe ser estacionario de la clase seleccionada.
Debe de ternerse en cuenta que para todas las variables declaradas en ambos niveles debe de definirse el valor de la componente.

\begin{lstlisting}
    SetPoint(model,Dict(x_1=>1,x_2=>1,y_1=>1,y_2=>1))
\end{lstlisting}

\subsubsection{Generar el Problema}
Finalmente al tener todos los datos previos introducidos se llama al \textbf{CreateProblem} pasandole como parámetro el \textit{model} 
y generará dicho problema imprimiendo en consola este.

\begin{lstlisting}
    CreateProblem(model)
\end{lstlisting}


\subsubsection{Documentación Oficial}
Para mayor información visitar la \href{https://github.com/FVSB/Tesis}{documentación oficial del generador}. 

Se pone antes de entrega rel host donde esta la pagina de MKDOCS
%TODO: Referenciar al link de usuario
	%
	%\include{Experimentación.tex}
	%
	%\chapter{Conclusiones}

%Explicar que se hizo
En esta tesis se desarrolló un algoritmo para la generación de problemas de optimización binivel con características específicas de estacionariedad en puntos determinados. El mismo tiene la capacidad de modificar problemas originales para garantizar la factibilidad y estacionariedad en puntos dados, aprovechando las capacidades del lenguaje de programación Julia, que destaca por su alto rendimiento computacional y sintaxis adaptada a problemas de optimización.

% Explicar como fue la experimentacion
La experimentación se llevó a cabo sobre tres categorías fundamentales de problemas: lineales, cuadráticos y no convexos. El proceso experimental comenzó con la obtención de puntos mínimos utilizando bibliotecas establecidas como BilevelJuMP y JuMP. Posteriormente, se generaron problemas estacionarios mediante la adición de componentes aleatorias a estos puntos para cada clase de problema definida.

% Que se comparó
Los problemas modificados fueron sometidos nuevamente a algoritmos tradicionales implementados en Julia para evaluar su efectividad frente a los puntos estacionarios generados. Se realizó un análisis comparativo destacando los casos más relevantes de cada categoría de puntos estacionarios, considerando la evaluación de la función objetivo del nivel superior en el punto donde se garantizó la estacionariedad, contrastándola con los resultados obtenidos por las bibliotecas convencionales.

% Recomendaciones
Como resultado de la investigación realizada, se han identificado varias líneas de trabajo futuro que permitirían expandir y mejorar los resultados obtenidos. En primer lugar, se recomienda ampliar el alcance de la experimentación numérica para incluir una mayor diversidad de problemas de optimización binivel. Esta expansión permitiría validar la robustez y versatilidad del algoritmo propuesto en diferentes contextos y escenarios de aplicación.

En segunda instancia, se sugiere profundizar en la investigación sobre la implementación del algoritmo desarrollado como criterio de parada en nuevos métodos de optimización binivel. Esta línea de investigación podría contribuir significativamente al desarrollo de algoritmos más eficientes y confiables, mejorando la capacidad de detectar y verificar puntos estacionarios durante el proceso de optimización.

Finalmente, se propone el desarrollo de una interfaz gráfica más intuitiva y funcional que facilite la generación automática de puntos según el tipo de estacionariedad requerida. Esta mejora en la usabilidad del software permitiría que usuarios con diferentes niveles de experiencia puedan aprovechar las capacidades del algoritmo de manera más efectiva, ampliando así su aplicabilidad práctica en diversos campos de estudio.
	%
	\backmatter
	\printbibliography

\end{document}